\chapter{Supervision and alerting}
\label{chap:supervision}
% Corresponds to phase 03 of the roadmap.

\section{The DevOps Control Center dashboard}
The application used throughout this internship to validate every mechanism described so far is itself a purpose-built supervision tool, developed from scratch as part of the internship: a Node.js/Express backend with a dependency-free vanilla JavaScript frontend (no build step), organised around seven views -- overview, applications, pipelines, resources, security, server, and settings -- with light/dark themes and a French/English interface. Its visual design was deliberately revisited midway through the internship, moving from a generic look to a dedicated design system (dense information layout, a consistent colour vocabulary for status, monospaced figures for machine-readable values) more appropriate for an operations tool that is scanned rather than read.

\section{SSH monitoring}
A background task tails the VM's \texttt{auth.log} continuously and classifies every connection attempt (accepted, failed, invalid user), enriching each one with an approximate geolocation of the source IP obtained from a free public API and cached locally so the same address is never looked up twice. This feed is what chapter~\ref{chap:securisation}'s brute-force and SSH-hardening work was validated against in practice.

\section{Supervision of the physical host (Proxmox)}
Beyond the guest VM, the dashboard also queries the Proxmox hypervisor itself through a read-only API token, surfacing the state and allocated resources of every virtual machine or container it hosts, along with basic disk health -- giving a single place to check both the application layer and the physical layer it runs on.

Two details deserve to be recorded because they cost time. The Proxmox log filter expects the name of the systemd unit (\texttt{ssh}) and not the name of the process (\texttt{sshd}); with the wrong name the API silently returns an empty result instead of an error. And once password login is disabled, the SSH daemon no longer writes ``Failed password'' lines for a valid user, so the log parser had to recognise three other wordings (connection closed, disconnected, connection reset by an authenticating user), otherwise only successful logins would have appeared and the history would have given a false impression of calm. The same parser feeds a filter in the security tab that shows the attempts on the virtual machine or on the physical server.

\section{Automatic verification of the server checklist}
The server tab holds a checklist of the hardening measures. Rather than relying on ticking boxes by hand, nine of its items are verified automatically by reading the real configuration on the host through read-only mounts: Docker installed, reverse proxy running, application deployed, pipelines configured, static IP, key-only SSH, active firewall, active Fail2ban and scheduled backups. A button next to each item triggers the check and shows the result; items that need an outside viewpoint (domain name, HTTPS certificate, access from outside the network) remain manual by design.

\section{Disk usage and the last Jenkins build}
Two additions came from real incidents. The resources tab breaks down the space used by folder on the whole host, the equivalent of \texttt{du} on the root directory; it was added after GitLab's disk filled up and is what allowed the cause to be found without an SSH session (chapter~\ref{chap:multiapps}). The pipelines tab also shows the status of the last Jenkins build (success, failure, unstable, running). Getting that status exposed a Docker networking limitation: the dashboard container timed out when calling Jenkins on the address of its own host machine, whatever the firewall rules, because of a ``hairpin'' problem. The solution was to give the container an alias (\texttt{host.docker.internal}) that points to the host, and to allow in the firewall the subnet that Docker actually uses for the project, not the default bridge.

\section{Proactive alerts}
Rather than requiring someone to have the dashboard open to notice a problem, two background watchers push alerts to a Discord channel through a webhook: one on container state transitions (a container that stops, crashes or comes back, checked every thirty seconds), and one on suspicious SSH activity (checked every fifteen seconds, with at most one message every five minutes so that an attack does not flood the channel). The backup check also sends its own alert. Discord webhooks were chosen over email specifically because they require no SMTP configuration and no account beyond a channel to alert into, which matched the internship's zero-budget constraint. Both watchers deliberately stay silent on their very first scan after startup, so that pre-existing history (a log already containing old entries, containers already running) is never mistaken for a fresh event.

\clearpage
